% ltex: language=it
\documentclass{article}

% Shared preamble from project root
\usepackage{lecture}

\begin{document}

\section{Contratto per la Manutenzione ed Evoluzione dei SI}\label{sec:contratto_per_la_manutenzione_ed_evoluzione_dei_si} % (fold)
Contratto sviluppato da CONSIP - società pubblica che fornisce consulenza tecnologica, organizzativa e di progetto agli uffici del Ministero dell'Economia.
\emptyln{}
Il contratto considerato prevede servizi di:
\begin{itemize}
    \item sviluppo e manutenzione evolutiva di software ad hoc;
    \item gestione applicativa;
    \item manutenzione adeguativa e correttiva;
    \item supporto specialistico;
\end{itemize}
\subsection{Parte Generale e Parte Speciale}\label{sub:parte_generale_e_parte_speciale} % (fold)
Molti enti utilizzano una struttura in due parti:
\begin{itemize}
    \item \textbf{generale}, comune a tutti i contratti, con gli aspetti standard;
    \item \textbf{speciale}, diversa di volta in volta, con gli aspetti peculiari;
\end{itemize}
%
\begin{multicols}{2}
\subsubsection{Parte Generale}\label{sub:parte_generale} % (fold)
    \begin{itemize}
        \item Aumento e diminuzione (standard 20\%, ``sesto quinto'');
        \item Modalità di esecuzione: luogo e convivenza, impiego di risorse specializzate;
        \item Rispetto della normativa sui rapporti di lavoro con possibile penalizzazione;
        \item Obblighi di riservatezza, con possibile risoluzione;
        \item Brevetti e diritti d'autore, rispetto garantito dal fornitore;
        \item Utilizzo di hw e sw: l'impresa deve essere autorizzata;
        \item Danni, responsabilità civile e assicurazione;
        \item Oneri fiscali e spese contrattuali;
        \item Cauzione: il 10\% del valore contrattuale, ridotto se l'azienda è certificata per la qualità, aumentato in caso di ribasso significativo; il costo di una polizza per la cauzione è 0,5-1\% del valore assicurato;
        \item Recesso per giusta causa: il committente deve dare un preavviso ma il fornitore no;
        \item Divieto di cessione del contratto e di cessione del credito;
        \item Trasparenza dei prezzi: assenza di intermediazione, rispetto della concorrenza;
        \item Subappalto: previsto a priori e con vincoli e responsabilità che restano sul fornitore e si ripetono sul subappaltatore;
        \item Foro competente esclusivo;
        \item Trattamento dati personali;
        \item Condizioni particolari di risoluzione;
    \end{itemize}
% subsubsection Parte Generale (end)
%
\subsubsection{Parte Speciale}\label{sub:parte_speciale} % (fold)
    \begin{itemize}
        \item Oggetto, luogo della prestazione e responsabile del proc.;
        \item Durata e affiancamento;
        \item Obblighi e adempimenti a carico dell'impresa;
        \item Proprietà del sw sviluppato e dei prodotti in genere;
        \item Dimensioni massime dei singoli servizi;
        \item Piano della qualità;
        \item Garanzie;
        \item Subappalto;
        \item Pianificazione delle attività;
        \item Produttività e risorse impiegate;
        \item Consegna dei prodotti;
        \item Collaudo e accettazione;
        \item Monitoraggio;
        \item Penali;
        \item Corrispettivo;
        \item Fatturazione;
        \item Risoluzione;
    \end{itemize}
% subsubsection Parte Speciale (end)
\end{multicols}
% subsection Parte Generale e Parte Speciale (end)
%
\subsection{Oggetto, luogo della prestazione e responsabile del proc}\label{sub:oggetto_luogo_della_prestazione_e_responsabile_del_proc} % (fold)
Servizi di:
\begin{itemize}
    \item sviluppo e manutenzione evolutiva di software ad hoc;
    \item gestione applicativa;
    \item manutenzione adeguativa e correttiva;
    \item supporto specialistico;
\end{itemize}
% subsection Oggetto, luogo della prestazione e responsabile del proc (end)
%
\subsection{Durata Affiancamento}\label{sub:durata_affiancamento} % (fold)
60 mesi di cui gli ultimi 12 solo per garanzia.
Nei primi due, affiancamento del fornitore uscente e negli ultimi due trasferimento di know-how al committente o a terzi.
% subsection Durata Affiancamento (end)
%
\subsection{Obblighi e Adempimenti a carico dell'impresa}\label{sub:obblighi_e_adempimenti_a_carico_dell_impresa} % (fold)
\begin{itemize}
    \item Oneri a rischi a carico dell'impresa, inclusi viaggi e missioni se necessari;
    \item Esecuzione a regola d'arte e con rispetto di regole tecniche e norme di sicurezza;
    \item Rispetto delle indicazioni del committente;
    \item Rispetto dei requisiti di accessibilità Web;
    \item Disponibilità a verifiche da parte del committente;
    \item Risoluzione e danni;
\end{itemize}
% subsection Obbligi e Adempimenti a carico dell'impresa (end)
%
\subsection{Proprietà del sw sviluppato}\label{sub:propriet_del_sw_sviluppato} % (fold)
L'amministrazione acquisisce la proprietà, con diritto di sfruttamento, di sw e della documentazione.
Essa può acquistare le licenze dei pacchetti utilizzati dal fornitore, oppure di utilizzare componenti open source.
% subsection Proprietà del sw sviluppato (end)
%
\subsection{Dimensioni massime dei singoli servizi}\label{sub:dimensioni_massime_dei_singoli_servizi} % (fold)
\begin{description}
    \item[Sviluppo e Manutenzione evolutiva] $60200$ punti funzione ($59800$ nuovi e $4000$ eliminati, pesati al $10\%$);
    \item[Gestione Applicativa] $38320$ giorni persona;
    \item[Manutenzione Adeguativa] $1840$ giorni persona;
    \item[Manutenzione Correttiva] con riferimento a $210900$ punti funzione;
    \item[Supporto Specialistico] $6520$ giorni persona con possibilità di travaso;
\end{description}
% subsection Dimensioni massime dei singoli servizi (end)
%
\subsection{Piano della qualità}\label{sub:piano_della_qualit} % (fold)
Il fornitore deve predisporre:
\begin{itemize}
    \item il \textbf{Piano della Qualità} generale;
    \item i \textbf{Piani della Qualità} per i vari obiettivi;
\end{itemize}
\begin{note}
    Il committente verifica e può richiedere modifiche
\end{note}
Per la struttura si utilizza: \textbf{Appendice 6 del Capitolato, par 2.1}.
\begin{itemize}
    \item Descrive un servizio in termini di:
        \begin{itemize}
            \item struttura organizzativa;
            \item responsabilità e risorse impiegate;
            \item procedure e procedimenti;
            \item spiega ``CHI'' fa ``COSA'', ``COME'' la fa e ``QUANDO'';
        \end{itemize}
    \item Assicura la qualità del servizio: controllo di processo e descrizione dei metodi di lavoro;
    \item Garantisce il fornitore che lo usa e il cliente che usufruisce dei servizi sviluppati;
\end{itemize}
%
\subsubsection{Indicatori di Qualità}\label{sub:indicatori_di_qualit} % (fold)
\begin{cimage}
    \includegraphics[width=.8\textwidth]{assets/indic-qual-es.png}
\end{cimage}
% subsubsection Indicatori di Qualità (end)
%
\subsubsection{Livelli di Servizio}\label{sub:livelli_di_servizio} % (fold)
Sono \textbf{elementi quantitativi} volti a definire soglie minime di accettazione per i vari elementi della fornitura: misurano il \textbf{valore rappresentato} da un \textbf{attributo di servizio}.
Due punti di vista:
\begin{itemize}
    \item tecnico (livelli di servizio):
        \begin{itemize}
            \item disponibilità del server;
            \item disponibilità della rete;
            \item disponibilità delle postazioni di lavoro;
        \end{itemize}
    \item utente (requisiti di servizio):
        \begin{itemize}
            \item disponibilità complessiva del sito Web;
            \item tempo massimo di interruzione;
        \end{itemize}
\end{itemize}
% subsubsection Livelli di Servizio (end)
%
\subsubsection{Aspetti tecnici e utente}\label{sub:aspetti_tecnici_e_utente} % (fold)
Fase di stesura del contratto (negoziazione):
\begin{itemize}
    \item l'\textbf{utente} esprime esigenze formalizzate nei requisiti da raggiungere a partire dalla situazione attuale;
    \item la \textbf{funzione informatica} rappresenta i requisiti degli utenti per il tramite di livelli di servizio;
    \item bisogna correlare le esigenze di servizio degli utenti (\textbf{Requisiti di Servizio}) e i criteri tecnici (\textbf{Livelli di Servizio});
\end{itemize}
% subsubsection Aspetti tecnici e utente (end)
% subsection Piano della qualità (end)
%
\subsection{Garanzie}\label{sub:garanzie} % (fold)
Il fornitore garantisce l'eliminazione di eventuali difetti riscontrati, con opportune modalità nell'ultimo anno.
% subsection Garanzie (end)
%
\subsection{Subappalto}\label{sub:subappalto} % (fold)
Vengono indicate:
\begin{itemize}
    \item prestazioni subappaltate;
    \item subappaltatori;
\end{itemize}
% subsection Subappalto (end)
% section Contratto per la Manutenzione ed Evoluzione dei SI (end)
\end{document}
